\section{Objectives}
\label{sec:objectives}

The project has an information-processing half and a networking half. The
information-processing half is the phrase on the phone and the check on the
laptop. The networking half is the decision of when that check is allowed
to occupy the uplink. Table~\ref{tab:objectives} lists the objectives and
the midterm status of each one.

\begin{table}[t]
\caption{Objectives and midterm status}
\label{tab:objectives}
\centering
\footnotesize
\begin{tabular}{p{0.72\columnwidth}c}
\toprule
Objective & Status \\
\midrule
Write an English action phrase on the phone before any picture is uploaded &
Done \\
Keep the picture on the phone when the phrase is confident and still matches the ongoing action &
Done \\
Let the laptop write its own phrase and compare the two phrases by text cosine &
Done \\
Check a short causal clip of the last three frames when the local test fails &
Done \\
Score free phrases on 60 still images and on three-frame UCF101 and HMDB51 clips &
Done \\
Predict checker latency, set a per-frame uplink budget, and schedule several phones &
Planned \\
\bottomrule
\end{tabular}
\end{table}

The first objective is a phrase, not a class identifier. The phone model
must finish generation, and the page must show the phrase, without waiting
for the laptop. The confidence attached to that phrase is the minimum
probability of the new tokens, so a single uncertain token can force a
check.

The second objective is the upload gate. A high confidence is not enough on
its own once an action is already on screen. The new phrase must also stay
close to that action. The picture is the payload that this gate can
suppress. The phrase itself may still be sent, because the text encoder
runs on the laptop.

The third objective replaces the proposal's bank caption. The laptop model
writes a free phrase. Agreement is the cosine of the two phrase embeddings.
The screen keeps the model's words. A class sentence is introduced only in
the evaluation, where a free phrase has to be counted against a label.

The fourth objective keeps the check causal and short. The phone reads the
newest frame. The laptop, and only the laptop, may read the last three
frames, captured about one second apart and uploaded in time order. The
system does not wait for a future frame, and it does not claim a
ten-second clip.

The fifth objective is a measurement that can be repeated. The models are
not trained on either sample. The still-image check scores one free phrase
against 15 class sentences. The video check replays the live rule on three
frames from each clip in the UCF101 recognition test list and the HMDB51
test list. Those frames are taken around the middle of the clip. The
replay therefore scores the gate, the quiet skip, and the phrase, and it
does not score the start or the end of an action.

The sixth objective is unchanged from the proposal and is not met at this
midterm. Before an upload, the laptop should predict how long the checker
will take, subtract that time from the confirmation deadline, and turn the
remainder into an uplink budget. When several phones ask for a check at
once, the earliest uplink deadline should go first. The November work is
to collect the traces, train the predictor, and compare that rule with
round-robin and with a static budget.
